\documentclass[11pt,a4paper]{article}
\newcommand{\gradelabel}{Grade 8}
\newcommand{\leconhdr}{Défi parcours robotisé}
% Préambule commun aux plans de cours de robotique (Grades 7 à 12).
% Chaque plan de cours commence par \documentclass[11pt,a4paper]{article}, définit
% \gradelabel et \leconhdr (en-têtes de page), puis % Préambule commun aux plans de cours de robotique (Grades 7 à 12).
% Chaque plan de cours commence par \documentclass[11pt,a4paper]{article}, définit
% \gradelabel et \leconhdr (en-têtes de page), puis % Préambule commun aux plans de cours de robotique (Grades 7 à 12).
% Chaque plan de cours commence par \documentclass[11pt,a4paper]{article}, définit
% \gradelabel et \leconhdr (en-têtes de page), puis \input{../../preambule_robotique.tex}.
\usepackage[T1]{fontenc}
\usepackage[utf8]{inputenc}
\usepackage{lmodern}
\usepackage[french]{babel}
\usepackage[margin=2cm,headheight=14pt]{geometry}
\usepackage[most]{tcolorbox}
\usepackage{xcolor,tabularx,booktabs,array,enumitem,fancyhdr,amssymb,tikz}
\usetikzlibrary{arrows.meta,positioning}

\pagestyle{fancy}
\fancyhf{}
\lhead{\small Robotique -- \gradelabel{} -- Novembre 2026}
\rhead{\small \leconhdr}
\cfoot{\thepage}
\setlist{itemsep=1pt,topsep=2pt,leftmargin=*}
\renewcommand{\arraystretch}{1.35}
\newcolumntype{Y}{>{\raggedright\arraybackslash}X}
\newcommand{\cb}{$\square$\,}
\newcommand{\fl}{$\rightarrow$}
\newcommand{\lig}{\noindent\rule{\linewidth}{0.3pt}\\[5mm]}

\newtcolorbox{etape}[3]{enhanced,breakable,colback=#1!6,colframe=#1!70!black,
  coltitle=white,fonttitle=\bfseries,title={#2\hfill #3},boxrule=0.8pt,arc=2mm,
  before skip=8pt,after skip=8pt}
\newtcolorbox{astuce}[1][Différenciation et besoins particuliers]{breakable,
  colback=yellow!12,colframe=yellow!50!black,fonttitle=\bfseries\small,
  title={#1},boxrule=0.5pt,arc=1mm,before skip=6pt,after skip=6pt}
\newtcolorbox{code}[1][Code]{breakable,colback=black!4,colframe=black!60,
  fonttitle=\bfseries\small,title={#1},boxrule=0.5pt,arc=1mm,fontupper=\ttfamily\small,
  before skip=6pt,after skip=6pt}
% Blocs de code coloré (écrire le code sans accents : comme dans l'éditeur)
\newtcblisting{arduino}[1][Code Arduino]{listing only,breakable,colback=black!4,
  colframe=black!60,fonttitle=\bfseries\small,title={#1},boxrule=0.5pt,arc=1mm,
  before skip=6pt,after skip=6pt,
  listing options={language=C++,basicstyle=\ttfamily\small,columns=fullflexible,
  keepspaces=true,showstringspaces=false,commentstyle=\color{green!40!black}\itshape,
  keywordstyle=\color{blue!70!black}\bfseries}}
\newtcblisting{python}[1][Code Python]{listing only,breakable,colback=black!4,
  colframe=black!60,fonttitle=\bfseries\small,title={#1},boxrule=0.5pt,arc=1mm,
  before skip=6pt,after skip=6pt,
  listing options={language=Python,basicstyle=\ttfamily\small,columns=fullflexible,
  keepspaces=true,showstringspaces=false,commentstyle=\color{green!40!black}\itshape,
  keywordstyle=\color{blue!70!black}\bfseries}}
\newcommand{\entete}[2]{%
  \begin{tcolorbox}[colback=#1!10,colframe=#1!70!black,boxrule=0.8pt,arc=2mm]
  {\Large\bfseries #2}\hfill Nom : \rule{4cm}{0.3pt}\quad Date : \rule{2.2cm}{0.3pt}
  \end{tcolorbox}}

% Bandeau de titre : #1 titre, #2 semaine
\newcommand{\titreblock}[2]{%
\begin{tcolorbox}[colback=teal!10,colframe=teal!70!black,arc=2mm,boxrule=1pt]
\centering
{\LARGE\bfseries #1}\\[2mm]
{\large Robotique -- \gradelabel{} -- Plan de cours de 2 heures}\\[1mm]
Enseignant : Francis Ki \quad|\quad #2
\end{tcolorbox}}

% Tableau du minutage : 6 contenus (accroche, mini-leçon, guidée, autonome, mise en commun, exit ticket)
\newcommand{\minutage}[6]{%
\begin{tabularx}{\linewidth}{>{\bfseries}p{2.2cm}p{2.6cm}Y p{3cm}}
\toprule
Minutes & Étape & Contenu & Organisation\\
\midrule
0 -- 10 & Accroche & #1 & Classe entière\\
10 -- 30 & Mini-leçon & #2 & Classe entière\\
30 -- 60 & Pratique guidée & #3 & Collectif puis binômes\\
60 -- 95 & Pratique autonome & #4 & Individuel sur poste\\
95 -- 110 & Mise en commun & #5 & Classe entière\\
110 -- 120 & Exit ticket & #6 & Individuel sur poste\\
\bottomrule
\end{tabularx}}

% Étayage progressif : #1 aide 1, #2 aide 2, #3 aide 3
\newcommand{\aides}[3]{%
\textbf{Étayage progressif de l'aide.}
\begin{center}
\begin{tabularx}{\linewidth}{>{\bfseries}p{2.4cm}Y p{4.4cm}}
\toprule
Niveau d'aide & Contenu & Qui intervient\\
\midrule
Aide 1 & #1 & Le binôme seul\\
Aide 2 & #2 & Le binôme avec la carte d'aide\\
Aide 3 & #3 & L'enseignant\\
\bottomrule
\end{tabularx}
\end{center}}

% Les trois groupes de la pratique autonome : #1 tâche A, #2 tâche B, #3 tâche C
\newcommand{\groupes}[3]{%
\textbf{Consigne commune\,:} chaque élève ouvre le fichier de son groupe, le renomme avec son prénom et travaille seul sur son poste. L'enseignant passe d'abord auprès du groupe A pendant 10 min.

\medskip
\begin{tabularx}{\linewidth}{>{\bfseries}p{3cm}Y p{2.2cm}}
\toprule
Groupe & Tâche sur poste & Annexe\\
\midrule
A : en dessous du niveau & #1 & Annexe C\\
B : au niveau & #2 & Annexe D\\
C : au-dessus du niveau & #3 & Annexe E\\
\bottomrule
\end{tabularx}}

% Mise en commun : #1 groupe A, #2 groupe B, #3 groupe C, #4 conclusion
\newcommand{\miseencommun}[4]{%
\begin{etape}{teal}{5. Mise en commun}{95 -- 110 min}
\begin{enumerate}
 \item (2 min) Retour au calme, écrans mis en veille.
 \item (9 min) Trois passages courts de 3 min, avec valorisation :
 \begin{itemize}
  \item \textbf{Groupe A} : #1
  \item \textbf{Groupe B} : #2
  \item \textbf{Groupe C} : #3
 \end{itemize}
 \item (4 min) Conclusion : #4
\end{enumerate}
\end{etape}}

% Clôture : #1 question 1, #2 question 2
\newcommand{\cloture}[2]{%
\begin{etape}{gray}{6. Clôture : exit ticket}{110 -- 120 min}
Formulaire en ligne ou papier (annexe G) :
\begin{enumerate}
 \item #1
 \item #2
 \item Auto-évaluation : \emph{Je maîtrise / Je progresse / J'ai besoin d'aide.}
\end{enumerate}
Lire les réponses avant la séance suivante pour ajuster les groupes de besoin.
\end{etape}}

% Grille d'observation : #1 #2 #3 intitulés de colonnes
\newcommand{\grille}[3]{%
\subsection*{Grille d'observation (enseignant)}
\begin{tabularx}{\linewidth}{Y p{2.4cm} p{2.4cm} p{2.4cm} Y}
\toprule
Élève & #1 & #2 & #3 & Groupe suivant\\
\midrule
 & & & & \\ & & & & \\ & & & & \\ & & & & \\ & & & & \\
\bottomrule
\end{tabularx}}

% Signature de l'exit ticket
\newcommand{\autoeval}{Je me sens : \cb Je maîtrise \quad \cb Je progresse \quad \cb J'ai besoin d'aide}
.
\usepackage[T1]{fontenc}
\usepackage[utf8]{inputenc}
\usepackage{lmodern}
\usepackage[french]{babel}
\usepackage[margin=2cm,headheight=14pt]{geometry}
\usepackage[most]{tcolorbox}
\usepackage{xcolor,tabularx,booktabs,array,enumitem,fancyhdr,amssymb,tikz}
\usetikzlibrary{arrows.meta,positioning}

\pagestyle{fancy}
\fancyhf{}
\lhead{\small Robotique -- \gradelabel{} -- Novembre 2026}
\rhead{\small \leconhdr}
\cfoot{\thepage}
\setlist{itemsep=1pt,topsep=2pt,leftmargin=*}
\renewcommand{\arraystretch}{1.35}
\newcolumntype{Y}{>{\raggedright\arraybackslash}X}
\newcommand{\cb}{$\square$\,}
\newcommand{\fl}{$\rightarrow$}
\newcommand{\lig}{\noindent\rule{\linewidth}{0.3pt}\\[5mm]}

\newtcolorbox{etape}[3]{enhanced,breakable,colback=#1!6,colframe=#1!70!black,
  coltitle=white,fonttitle=\bfseries,title={#2\hfill #3},boxrule=0.8pt,arc=2mm,
  before skip=8pt,after skip=8pt}
\newtcolorbox{astuce}[1][Différenciation et besoins particuliers]{breakable,
  colback=yellow!12,colframe=yellow!50!black,fonttitle=\bfseries\small,
  title={#1},boxrule=0.5pt,arc=1mm,before skip=6pt,after skip=6pt}
\newtcolorbox{code}[1][Code]{breakable,colback=black!4,colframe=black!60,
  fonttitle=\bfseries\small,title={#1},boxrule=0.5pt,arc=1mm,fontupper=\ttfamily\small,
  before skip=6pt,after skip=6pt}
% Blocs de code coloré (écrire le code sans accents : comme dans l'éditeur)
\newtcblisting{arduino}[1][Code Arduino]{listing only,breakable,colback=black!4,
  colframe=black!60,fonttitle=\bfseries\small,title={#1},boxrule=0.5pt,arc=1mm,
  before skip=6pt,after skip=6pt,
  listing options={language=C++,basicstyle=\ttfamily\small,columns=fullflexible,
  keepspaces=true,showstringspaces=false,commentstyle=\color{green!40!black}\itshape,
  keywordstyle=\color{blue!70!black}\bfseries}}
\newtcblisting{python}[1][Code Python]{listing only,breakable,colback=black!4,
  colframe=black!60,fonttitle=\bfseries\small,title={#1},boxrule=0.5pt,arc=1mm,
  before skip=6pt,after skip=6pt,
  listing options={language=Python,basicstyle=\ttfamily\small,columns=fullflexible,
  keepspaces=true,showstringspaces=false,commentstyle=\color{green!40!black}\itshape,
  keywordstyle=\color{blue!70!black}\bfseries}}
\newcommand{\entete}[2]{%
  \begin{tcolorbox}[colback=#1!10,colframe=#1!70!black,boxrule=0.8pt,arc=2mm]
  {\Large\bfseries #2}\hfill Nom : \rule{4cm}{0.3pt}\quad Date : \rule{2.2cm}{0.3pt}
  \end{tcolorbox}}

% Bandeau de titre : #1 titre, #2 semaine
\newcommand{\titreblock}[2]{%
\begin{tcolorbox}[colback=teal!10,colframe=teal!70!black,arc=2mm,boxrule=1pt]
\centering
{\LARGE\bfseries #1}\\[2mm]
{\large Robotique -- \gradelabel{} -- Plan de cours de 2 heures}\\[1mm]
Enseignant : Francis Ki \quad|\quad #2
\end{tcolorbox}}

% Tableau du minutage : 6 contenus (accroche, mini-leçon, guidée, autonome, mise en commun, exit ticket)
\newcommand{\minutage}[6]{%
\begin{tabularx}{\linewidth}{>{\bfseries}p{2.2cm}p{2.6cm}Y p{3cm}}
\toprule
Minutes & Étape & Contenu & Organisation\\
\midrule
0 -- 10 & Accroche & #1 & Classe entière\\
10 -- 30 & Mini-leçon & #2 & Classe entière\\
30 -- 60 & Pratique guidée & #3 & Collectif puis binômes\\
60 -- 95 & Pratique autonome & #4 & Individuel sur poste\\
95 -- 110 & Mise en commun & #5 & Classe entière\\
110 -- 120 & Exit ticket & #6 & Individuel sur poste\\
\bottomrule
\end{tabularx}}

% Étayage progressif : #1 aide 1, #2 aide 2, #3 aide 3
\newcommand{\aides}[3]{%
\textbf{Étayage progressif de l'aide.}
\begin{center}
\begin{tabularx}{\linewidth}{>{\bfseries}p{2.4cm}Y p{4.4cm}}
\toprule
Niveau d'aide & Contenu & Qui intervient\\
\midrule
Aide 1 & #1 & Le binôme seul\\
Aide 2 & #2 & Le binôme avec la carte d'aide\\
Aide 3 & #3 & L'enseignant\\
\bottomrule
\end{tabularx}
\end{center}}

% Les trois groupes de la pratique autonome : #1 tâche A, #2 tâche B, #3 tâche C
\newcommand{\groupes}[3]{%
\textbf{Consigne commune\,:} chaque élève ouvre le fichier de son groupe, le renomme avec son prénom et travaille seul sur son poste. L'enseignant passe d'abord auprès du groupe A pendant 10 min.

\medskip
\begin{tabularx}{\linewidth}{>{\bfseries}p{3cm}Y p{2.2cm}}
\toprule
Groupe & Tâche sur poste & Annexe\\
\midrule
A : en dessous du niveau & #1 & Annexe C\\
B : au niveau & #2 & Annexe D\\
C : au-dessus du niveau & #3 & Annexe E\\
\bottomrule
\end{tabularx}}

% Mise en commun : #1 groupe A, #2 groupe B, #3 groupe C, #4 conclusion
\newcommand{\miseencommun}[4]{%
\begin{etape}{teal}{5. Mise en commun}{95 -- 110 min}
\begin{enumerate}
 \item (2 min) Retour au calme, écrans mis en veille.
 \item (9 min) Trois passages courts de 3 min, avec valorisation :
 \begin{itemize}
  \item \textbf{Groupe A} : #1
  \item \textbf{Groupe B} : #2
  \item \textbf{Groupe C} : #3
 \end{itemize}
 \item (4 min) Conclusion : #4
\end{enumerate}
\end{etape}}

% Clôture : #1 question 1, #2 question 2
\newcommand{\cloture}[2]{%
\begin{etape}{gray}{6. Clôture : exit ticket}{110 -- 120 min}
Formulaire en ligne ou papier (annexe G) :
\begin{enumerate}
 \item #1
 \item #2
 \item Auto-évaluation : \emph{Je maîtrise / Je progresse / J'ai besoin d'aide.}
\end{enumerate}
Lire les réponses avant la séance suivante pour ajuster les groupes de besoin.
\end{etape}}

% Grille d'observation : #1 #2 #3 intitulés de colonnes
\newcommand{\grille}[3]{%
\subsection*{Grille d'observation (enseignant)}
\begin{tabularx}{\linewidth}{Y p{2.4cm} p{2.4cm} p{2.4cm} Y}
\toprule
Élève & #1 & #2 & #3 & Groupe suivant\\
\midrule
 & & & & \\ & & & & \\ & & & & \\ & & & & \\ & & & & \\
\bottomrule
\end{tabularx}}

% Signature de l'exit ticket
\newcommand{\autoeval}{Je me sens : \cb Je maîtrise \quad \cb Je progresse \quad \cb J'ai besoin d'aide}
.
\usepackage[T1]{fontenc}
\usepackage[utf8]{inputenc}
\usepackage{lmodern}
\usepackage[french]{babel}
\usepackage[margin=2cm,headheight=14pt]{geometry}
\usepackage[most]{tcolorbox}
\usepackage{xcolor,tabularx,booktabs,array,enumitem,fancyhdr,amssymb,tikz}
\usetikzlibrary{arrows.meta,positioning}

\pagestyle{fancy}
\fancyhf{}
\lhead{\small Robotique -- \gradelabel{} -- Novembre 2026}
\rhead{\small \leconhdr}
\cfoot{\thepage}
\setlist{itemsep=1pt,topsep=2pt,leftmargin=*}
\renewcommand{\arraystretch}{1.35}
\newcolumntype{Y}{>{\raggedright\arraybackslash}X}
\newcommand{\cb}{$\square$\,}
\newcommand{\fl}{$\rightarrow$}
\newcommand{\lig}{\noindent\rule{\linewidth}{0.3pt}\\[5mm]}

\newtcolorbox{etape}[3]{enhanced,breakable,colback=#1!6,colframe=#1!70!black,
  coltitle=white,fonttitle=\bfseries,title={#2\hfill #3},boxrule=0.8pt,arc=2mm,
  before skip=8pt,after skip=8pt}
\newtcolorbox{astuce}[1][Différenciation et besoins particuliers]{breakable,
  colback=yellow!12,colframe=yellow!50!black,fonttitle=\bfseries\small,
  title={#1},boxrule=0.5pt,arc=1mm,before skip=6pt,after skip=6pt}
\newtcolorbox{code}[1][Code]{breakable,colback=black!4,colframe=black!60,
  fonttitle=\bfseries\small,title={#1},boxrule=0.5pt,arc=1mm,fontupper=\ttfamily\small,
  before skip=6pt,after skip=6pt}
% Blocs de code coloré (écrire le code sans accents : comme dans l'éditeur)
\newtcblisting{arduino}[1][Code Arduino]{listing only,breakable,colback=black!4,
  colframe=black!60,fonttitle=\bfseries\small,title={#1},boxrule=0.5pt,arc=1mm,
  before skip=6pt,after skip=6pt,
  listing options={language=C++,basicstyle=\ttfamily\small,columns=fullflexible,
  keepspaces=true,showstringspaces=false,commentstyle=\color{green!40!black}\itshape,
  keywordstyle=\color{blue!70!black}\bfseries}}
\newtcblisting{python}[1][Code Python]{listing only,breakable,colback=black!4,
  colframe=black!60,fonttitle=\bfseries\small,title={#1},boxrule=0.5pt,arc=1mm,
  before skip=6pt,after skip=6pt,
  listing options={language=Python,basicstyle=\ttfamily\small,columns=fullflexible,
  keepspaces=true,showstringspaces=false,commentstyle=\color{green!40!black}\itshape,
  keywordstyle=\color{blue!70!black}\bfseries}}
\newcommand{\entete}[2]{%
  \begin{tcolorbox}[colback=#1!10,colframe=#1!70!black,boxrule=0.8pt,arc=2mm]
  {\Large\bfseries #2}\hfill Nom : \rule{4cm}{0.3pt}\quad Date : \rule{2.2cm}{0.3pt}
  \end{tcolorbox}}

% Bandeau de titre : #1 titre, #2 semaine
\newcommand{\titreblock}[2]{%
\begin{tcolorbox}[colback=teal!10,colframe=teal!70!black,arc=2mm,boxrule=1pt]
\centering
{\LARGE\bfseries #1}\\[2mm]
{\large Robotique -- \gradelabel{} -- Plan de cours de 2 heures}\\[1mm]
Enseignant : Francis Ki \quad|\quad #2
\end{tcolorbox}}

% Tableau du minutage : 6 contenus (accroche, mini-leçon, guidée, autonome, mise en commun, exit ticket)
\newcommand{\minutage}[6]{%
\begin{tabularx}{\linewidth}{>{\bfseries}p{2.2cm}p{2.6cm}Y p{3cm}}
\toprule
Minutes & Étape & Contenu & Organisation\\
\midrule
0 -- 10 & Accroche & #1 & Classe entière\\
10 -- 30 & Mini-leçon & #2 & Classe entière\\
30 -- 60 & Pratique guidée & #3 & Collectif puis binômes\\
60 -- 95 & Pratique autonome & #4 & Individuel sur poste\\
95 -- 110 & Mise en commun & #5 & Classe entière\\
110 -- 120 & Exit ticket & #6 & Individuel sur poste\\
\bottomrule
\end{tabularx}}

% Étayage progressif : #1 aide 1, #2 aide 2, #3 aide 3
\newcommand{\aides}[3]{%
\textbf{Étayage progressif de l'aide.}
\begin{center}
\begin{tabularx}{\linewidth}{>{\bfseries}p{2.4cm}Y p{4.4cm}}
\toprule
Niveau d'aide & Contenu & Qui intervient\\
\midrule
Aide 1 & #1 & Le binôme seul\\
Aide 2 & #2 & Le binôme avec la carte d'aide\\
Aide 3 & #3 & L'enseignant\\
\bottomrule
\end{tabularx}
\end{center}}

% Les trois groupes de la pratique autonome : #1 tâche A, #2 tâche B, #3 tâche C
\newcommand{\groupes}[3]{%
\textbf{Consigne commune\,:} chaque élève ouvre le fichier de son groupe, le renomme avec son prénom et travaille seul sur son poste. L'enseignant passe d'abord auprès du groupe A pendant 10 min.

\medskip
\begin{tabularx}{\linewidth}{>{\bfseries}p{3cm}Y p{2.2cm}}
\toprule
Groupe & Tâche sur poste & Annexe\\
\midrule
A : en dessous du niveau & #1 & Annexe C\\
B : au niveau & #2 & Annexe D\\
C : au-dessus du niveau & #3 & Annexe E\\
\bottomrule
\end{tabularx}}

% Mise en commun : #1 groupe A, #2 groupe B, #3 groupe C, #4 conclusion
\newcommand{\miseencommun}[4]{%
\begin{etape}{teal}{5. Mise en commun}{95 -- 110 min}
\begin{enumerate}
 \item (2 min) Retour au calme, écrans mis en veille.
 \item (9 min) Trois passages courts de 3 min, avec valorisation :
 \begin{itemize}
  \item \textbf{Groupe A} : #1
  \item \textbf{Groupe B} : #2
  \item \textbf{Groupe C} : #3
 \end{itemize}
 \item (4 min) Conclusion : #4
\end{enumerate}
\end{etape}}

% Clôture : #1 question 1, #2 question 2
\newcommand{\cloture}[2]{%
\begin{etape}{gray}{6. Clôture : exit ticket}{110 -- 120 min}
Formulaire en ligne ou papier (annexe G) :
\begin{enumerate}
 \item #1
 \item #2
 \item Auto-évaluation : \emph{Je maîtrise / Je progresse / J'ai besoin d'aide.}
\end{enumerate}
Lire les réponses avant la séance suivante pour ajuster les groupes de besoin.
\end{etape}}

% Grille d'observation : #1 #2 #3 intitulés de colonnes
\newcommand{\grille}[3]{%
\subsection*{Grille d'observation (enseignant)}
\begin{tabularx}{\linewidth}{Y p{2.4cm} p{2.4cm} p{2.4cm} Y}
\toprule
Élève & #1 & #2 & #3 & Groupe suivant\\
\midrule
 & & & & \\ & & & & \\ & & & & \\ & & & & \\ & & & & \\
\bottomrule
\end{tabularx}}

% Signature de l'exit ticket
\newcommand{\autoeval}{Je me sens : \cb Je maîtrise \quad \cb Je progresse \quad \cb J'ai besoin d'aide}


\begin{document}
\titreblock{Défi parcours robotisé}{Semaine du 23 au 27 novembre 2026}

%==================================================================
\section{Informations générales}
\begin{tabularx}{\linewidth}{>{\bfseries}p{3.6cm}Y}
\toprule
Titre & Défi parcours robotisé\\
Classe / effectif & Grade 8, 11 à 20 élèves, niveaux hétérogènes, élèves à besoins particuliers\\
Durée & 2 heures (120 min)\\
Lieu & Salle informatique, un poste par élève\\
Matériel et logiciels &
\begin{itemize}
 \item Un poste par élève avec mBlock 5, robot connecté.
 \item Un mBot par binôme, chargé.
 \item Un parcours au sol réalisé au ruban de masquage : départ, trois segments, deux virages à 90°, zone d'arrivée (plan en annexe A).
 \item Un chronomètre et un mètre ruban pour le défi.
\end{itemize}\\
Pré-requis &
Séances 1 à 3 : connexion, déplacements, rotations calibrées. Boucle \texttt{répéter} (G7).\\
\bottomrule
\end{tabularx}

\begin{astuce}[Préparation avant le cours (environ 45 min)]
\begin{enumerate}
 \item Tracer au moins deux exemplaires du parcours de l'annexe A (un pour les groupes A et B, un pour le groupe C), avec les distances écrites.
 \item Préparer les fichiers de départ et la grille d'évaluation du défi (annexe B).
 \item Vérifier la charge des batteries de tous les robots.
 \item Créer le formulaire de sortie (annexe G) ou imprimer l'annexe.
\end{enumerate}
\end{astuce}

%==================================================================
\section{Objectifs d'apprentissage}
\begin{tcolorbox}[colback=gray!8,colframe=gray!60!black,arc=1mm,boxrule=0.5pt]
\textbf{Objectifs du programme}
\begin{itemize}
 \item Réaliser un parcours imposé.
 \item Optimiser le programme du robot.
\end{itemize}
\textbf{À la fin de la séance, l'élève est capable\,:}
\begin{itemize}
 \item de décomposer un parcours en segments et en virages dans un tableau\,;
 \item de programmer, tester et corriger son programme jusqu'à atteindre la zone d'arrivée\,;
 \item d'optimiser un programme (moins de blocs avec des boucles ou des variables, plus de précision).
\end{itemize}
\textbf{Compétences numériques\,:} versionner (enregistrer chaque version), tester méthodiquement, comparer deux programmes.
\end{tcolorbox}

%==================================================================
\section{Déroulé minuté (120 min)}
\minutage{Présentation du défi et de ses règles}{Décomposer et optimiser un parcours}{Le tableau des segments en collectif, puis le premier segment en binômes}{Trois niveaux de besoin}{Passage des robots et valorisation}{Trois questions sur poste}

%------------------------------------------------------------------
\begin{etape}{orange}{1. Accroche : le défi}{0 -- 10 min}
\begin{enumerate}
 \item (4 min) Présenter le parcours au sol et ses règles : départ derrière la ligne, ne pas toucher le robot, arrivée dans la zone, pénalité si le robot sort.
 \item (3 min) Montrer les trois critères de réussite : \textbf{précision}, \textbf{optimisation} (peu de blocs), \textbf{persévérance} (nombre de corrections documentées).
 \item (3 min) Question-problème : \textbf{«~Comment programmer un parcours qu'on ne peut pas réussir du premier coup ?~»}
\end{enumerate}
\end{etape}

%------------------------------------------------------------------
\begin{etape}{blue}{2. Mini-leçon : décomposer et optimiser}{10 -- 30 min}
\textbf{Support\,:} annexe A. Quatre temps de 5 min.

\textbf{Temps 1 : décomposer.} Un parcours est une liste de segments et de virages. On la note dans un tableau.

\begin{center}
\begin{tabularx}{\linewidth}{>{\bfseries}p{1.6cm}Y Y Y}
\toprule
N° & Action & Mesure ou angle & Réglage à trouver\\
\midrule
1 & Avancer & 60 cm & durée\\
2 & Tourner & 90° à droite & durée\\
3 & Avancer & 40 cm & durée\\
\bottomrule
\end{tabularx}
\end{center}

\textbf{Temps 2 : programmer par étapes.} On programme un segment, on le teste, puis on ajoute le suivant. On n'écrit jamais tout d'un coup.

\textbf{Temps 3 : corriger.} Mesurer l'écart, ajuster un seul paramètre à la fois, noter le résultat.

\textbf{Temps 4 : optimiser.} Un programme est meilleur quand il est plus court (boucles), plus facile à modifier (variables comme \texttt{durée\_90}) et plus précis.

\begin{astuce}
Garder le tableau des segments affiché. Un seul réglage modifié à la fois. Fournir une checklist «~Programmer, Tester, Corriger~».
\end{astuce}
\end{etape}

%------------------------------------------------------------------
\begin{etape}{green}{3. Pratique guidée}{30 -- 60 min}
\textbf{Fiche\,:} annexe B.

\textbf{3a. Le tableau des segments en collectif (12 min).} Remplir ensemble le tableau avec les mesures du parcours projeté.

\textbf{3b. Le premier segment en binômes (18 min).} Chaque binôme programme et teste le segment 1 et le premier virage. Alternance des rôles.

\aides{Le tableau rempli du 3a et les durées calibrées aux séances précédentes}{Questions guides : \emph{Quelle est la mesure de ce segment ? Quelle durée avais-tu trouvée pour 90° ?}}{L'enseignant teste le premier segment avec le binôme, qui continue}

\begin{astuce}
Autoriser à réutiliser les valeurs calibrées des séances précédentes. Si un binôme a un robot qui dévie, lui proposer un robot de rechange.
\end{astuce}
\end{etape}

%------------------------------------------------------------------
\begin{etape}{violet}{4. Pratique autonome différenciée}{60 -- 95 min}
Chaque élève prépare son programme seul. Les essais se font par binôme sur les parcours, à tour de rôle.

\medskip
\groupes{Parcours simplifié à deux segments et un virage, avec un squelette de programme à compléter et des durées proposées. Mise en ordre des blocs avec des étiquettes de couleur.}{Programmer le parcours complet de l'annexe A, tester, corriger, et tenir un journal des corrections (une ligne par essai).}{Parcours complet avec optimisation : variables \texttt{duree\_90} et \texttt{duree\_60cm}, boucle \texttt{répéter}, retour au point de départ. Comparer le nombre de blocs avant et après optimisation.}

\medskip
\textbf{Plan B si un élève termine tôt\,:} réussir le parcours dans le temps le plus court possible sans sortir.
\end{etape}

%------------------------------------------------------------------
\begin{etape}{teal}{5. Mise en commun : passage des robots}{95 -- 110 min}
\begin{enumerate}
 \item (2 min) Retour au calme, écrans mis en veille.
 \item (9 min) Un passage officiel par binôme volontaire sur le parcours. La classe valorise un point précis :
 \begin{itemize}
  \item \textbf{Groupe A} : réussite du premier segment. Valoriser l'effort de réglage.
  \item \textbf{Groupe B} : réussite du parcours. Valoriser la précision et le journal des corrections.
  \item \textbf{Groupe C} : programme optimisé. Valoriser la réduction du nombre de blocs.
 \end{itemize}
 \item (4 min) Conclusion : \textit{un parcours se réussit en décomposant, en testant et en corrigeant, puis en optimisant.}
\end{enumerate}
\end{etape}

\cloture{Pourquoi programme-t-on un segment à la fois et non tout le parcours d'un coup ?}{Cite une façon d'optimiser un programme de robot.}

%==================================================================
% ANNEXES
%==================================================================
\newpage
\section*{Annexe A : Support résumé de la mini-leçon et plan du parcours}
\entete{blue}{Défi parcours robotisé}

\begin{center}
\begin{tikzpicture}[>={Stealth[length=3mm]}]
\draw[very thick,blue!60] (0,0)--(6,0)--(6,3)--(3,3);
\node[below] at (0,0) {Départ};
\node[above] at (3,3) {Arrivée};
\node[below] at (3,0) {60 cm};
\node[right] at (6,1.5) {40 cm};
\node[above] at (4.5,3) {30 cm};
\draw[->,thick] (0,0)--(0.6,0);
\node[circle,fill=orange!60,inner sep=2pt] at (6,0) {1};
\node[circle,fill=orange!60,inner sep=2pt] at (6,3) {2};
\end{tikzpicture}
\end{center}

\begin{tabularx}{\linewidth}{>{\bfseries}p{1.6cm}Y Y Y}
\toprule
N° & Action & Mesure ou angle & Réglage à trouver\\
\midrule
1 & Avancer & 60 cm & durée\\
2 & Tourner & 90° à gauche & durée\\
3 & Avancer & 40 cm & durée\\
4 & Tourner & 90° à gauche & durée\\
5 & Avancer & 30 cm & durée\\
\bottomrule
\end{tabularx}

\medskip
\textbf{Méthode :} Programmer un segment, tester, corriger (un réglage à la fois), noter. Optimiser à la fin.

%==================================================================
\newpage
\section*{Annexe B : Fiche de pratique guidée et grille du défi}
\entete{green}{Pratique guidée : mon parcours}

\subsection*{3a. Tableau des segments (à compléter avec la classe)}
\begin{tabularx}{\linewidth}{>{\bfseries}p{1.6cm}Y Y Y}
\toprule
N° & Action & Mesure ou angle & Durée trouvée (s)\\
\midrule
1 & & & \\ 2 & & & \\ 3 & & & \\ 4 & & & \\ 5 & & & \\
\bottomrule
\end{tabularx}

\subsection*{Grille du défi (enseignant et élèves)}
\begin{tabularx}{\linewidth}{>{\bfseries}p{3.6cm}Y p{2.4cm}}
\toprule
Critère & Description & Points\\
\midrule
Précision & Le robot atteint la zone d'arrivée sans sortir du parcours & /4\\
Optimisation & Programme avec boucle ou variable, peu de blocs & /3\\
Persévérance & Journal des corrections complet & /3\\
\bottomrule
\end{tabularx}

\begin{tcolorbox}[colback=yellow!10,colframe=yellow!50!black,title={Cartes d'aide à utiliser dans l'ordre},fonttitle=\bfseries\small]
\textbf{Aide 1 :} regarde le tableau rempli et les durées des séances précédentes.\\
\textbf{Aide 2 :} Quelle est la mesure de ce segment ? Quelle durée avais-tu trouvée pour 90° ?\\
\textbf{Aide 3 :} demande à l'enseignant de tester le premier segment avec toi.
\end{tcolorbox}

%==================================================================
\newpage
\section*{Annexe C : Niveau A (en dessous du niveau)}
\entete{orange}{Niveau A : parcours simplifié}

\textbf{Parcours :} avancer de 40 cm, tourner à droite de 90°, avancer de 30 cm.

\subsection*{Étape 1 : complète le programme}
\begin{enumerate}
 \item Avancer à la puissance 50 \% pendant \dotfill{} s.
 \item Arrêter les moteurs.
 \item Tourner à droite pendant \dotfill{} s.
 \item Arrêter les moteurs.
 \item Avancer pendant \dotfill{} s, puis arrêter.
\end{enumerate}

\subsection*{Étape 2 : journal des essais}
\begin{tabularx}{\linewidth}{>{\bfseries}p{2cm}Y Y}
\toprule
Essai & Ce que j'ai changé & Résultat\\
\midrule
1 & & \\ 2 & & \\ 3 & & \\
\bottomrule
\end{tabularx}

%==================================================================
\newpage
\section*{Annexe D : Niveau B (au niveau)}
\entete{blue}{Niveau B : le parcours complet}

\textbf{Mission :} réussir le parcours de l'annexe A.

\begin{tabularx}{\linewidth}{>{\bfseries}p{1.6cm}Y Y Y}
\toprule
Essai & Réglage modifié & Nouvelle valeur & Résultat\\
\midrule
1 & & & \\ 2 & & & \\ 3 & & & \\ 4 & & & \\ 5 & & & \\ 6 & & & \\
\bottomrule
\end{tabularx}

\subsection*{Bilan}
Le réglage le plus difficile à trouver était : \dotfill\\
Je l'ai trouvé en : \dotfill

%==================================================================
\newpage
\section*{Annexe E : Niveau C (au-dessus du niveau)}
\entete{violet}{Niveau C : optimiser et revenir}

\textbf{Mission :} réussir le parcours de l'annexe A, puis revenir au point de départ.

\begin{tabularx}{\linewidth}{>{\bfseries}p{4cm}Y Y}
\toprule
Version du programme & Nombre de blocs & Précision (cm d'écart)\\
\midrule
V1 (sans optimisation) & & \\
V2 (avec variables) & & \\
V3 (avec boucle) & & \\
\bottomrule
\end{tabularx}

\subsection*{Questions}
\begin{enumerate}
 \item Quelle optimisation a le plus réduit le nombre de blocs ? \dotfill
 \item Pourquoi une variable comme \texttt{duree\_90} facilite-t-elle les corrections ? \dotfill
\end{enumerate}

\textbf{Défi bonus :} réussir le parcours deux fois de suite sans modifier le programme.

%==================================================================
\newpage
\section*{Annexe F : Corrections et attendus (enseignant)}

\subsection*{Annexe B}
Le tableau doit contenir cinq lignes alternant avancer et tourner. Les durées dépendent du robot ; elles sont les mêmes que celles trouvées aux séances 2 et 3.

\subsection*{Annexe C}
Mêmes durées que dans les séances précédentes. Le journal doit montrer un changement à la fois.

\subsection*{Annexe D}
Le programme doit comporter cinq actions. Le réglage souvent le plus difficile est la rotation. Critères de réussite : arrivée dans la zone, journal avec au moins trois essais.

\subsection*{Annexe E}
Optimisation attendue : variables \texttt{duree\_90} et \texttt{duree\_cm}, ou boucle si les virages sont identiques. Les variables rendent les corrections globales en un seul endroit.\\
\textbf{Critères de réussite :} trois versions comparées, réduction du nombre de blocs.

%==================================================================
\newpage
\section*{Annexe G : Exit ticket et grille d'observation}
\entete{gray}{Exit ticket : défi parcours}

\begin{enumerate}
 \item Pourquoi programme-t-on un segment à la fois et non tout le parcours d'un coup ?\\ \lig
 \item Cite une façon d'optimiser un programme de robot.\\ \lig
 \item \autoeval
\end{enumerate}
\textbf{Attendus :} 1. Pour trouver plus facilement l'erreur et corriger un seul réglage. 2. Utiliser une boucle ou une variable, réduire le nombre de blocs.

\bigskip
\grille{Tableau des segments}{Journal des essais}{Optimisation}

\end{document}
